\documentclass[11pt]{article}
\usepackage[T1]{fontenc}
\usepackage{lmodern}
\usepackage[utf8]{inputenc}
\usepackage[english]{babel}
\usepackage{amsmath,amssymb,amsthm}
\usepackage{booktabs,array}
\usepackage[margin=1in]{geometry}
\usepackage{microtype}
\usepackage[colorlinks=true,linkcolor=black,citecolor=blue,urlcolor=blue]{hyperref}
\hypersetup{pdfauthor={Carlo Corti},pdftitle={Exact Enumeration of Four-Window Permutations: OEIS A180207 through n=22}}
\newtheorem{theorem}{Theorem}[section]
\newtheorem{proposition}[theorem]{Proposition}
\newtheorem{lemma}[theorem]{Lemma}
\newtheorem{corollary}[theorem]{Corollary}
\theoremstyle{definition}
\newtheorem{definition}[theorem]{Definition}
\theoremstyle{remark}
\newtheorem{remark}[theorem]{Remark}
\newcommand{\Znn}{\mathbb Z_{\geq0}}
\newcommand{\calA}{\mathcal A}
\title{Exact Enumeration of Four-Window Permutations:\\OEIS A180207 through $n=22$}
\author{Carlo Corti\\Independent researcher\\\texttt{carlo.corti@outlook.com}}
\date{November 2026}
\begin{document}
\maketitle

\begin{abstract}
OEIS A180207 counts ordered permutations of $0,\ldots,n-1$ that are
consecutive sums of four entries of a nonnegative integer sequence of length
$n+3$. We give an exact enumeration through $n=22$, extending the historical
16-term prefix used in this study by six locally computed values, ending in
$a(22)=2055385215352$. The zero entry forces four consecutive zeros in the
realization, makes that realization unique, and decomposes each permutation
into two outward tails. A weighted subset dynamic program counts these tails
and combines complementary label sets. A second complete count through $22$
uses inclusion--exclusion on walks with repeated labels; an unrestricted-window
Python method supplies an additional check at $1,\ldots,12,16,17$.
We prove a residue-minimum feasibility criterion, the counting identities,
and bounds for their intermediate arithmetic. Reversal gives evenness for
$n>1$, while two end-extension injections give $a(n+1)\geq2a(n)$.
\end{abstract}

\noindent\textbf{2020 Mathematics Subject Classification.}
Primary 05A05; Secondary 05A15, 05A19.

\smallskip
\noindent\textbf{Key words and phrases.}
permutation, consecutive sum, nonnegative realization, subset dynamic
programming, inclusion--exclusion, OEIS A180207.

\section{Introduction}
\label{sec:intro}

Consider a permutation $p=(p_0,\ldots,p_{n-1})$ of $\{0,\ldots,n-1\}$.
We ask whether there is an integer vector $x=(x_0,\ldots,x_{n+2})$ with
nonnegative entries such that
\begin{equation}
 p_i=x_i+x_{i+1}+x_{i+2}+x_{i+3}\qquad(0\leq i<n).
 \label{eq:windows}
\end{equation}
Let $\calA_n$ be the set of such permutations and let $a(n)=|\calA_n|$,
for $n\geq1$, with the same index $n$ and initial index $1$ as OEIS A180207.
The OEIS entry credits R. H. Hardin with author date
15 August 2010~\cite{hardin}. The OEIS snapshot used for the historical baseline
contains the values through $n=16$, ending in $35154332$. These are the
historical baseline of the present study and retain Hardin's attribution;
the entry's author date does not date the computation of each individual term.

Related OEIS entries replace the window length four by two, three, five,
six and seven: A180205, A180206, A180208, A180209 and A180210,
respectively~\cite{family}. These are variants of the definition; no exact
transform from their counts to $a(n)$ is used here. The forward problem of
counting distinct interval sums of a permutation is studied by
Konieczny~\cite{konieczny}, while Higashitani and Kurimoto~\cite{higashitani}
study the minimum possible maximal cyclic window sum. Their input is a
permutation, whereas here it is the output word of the window operation
that must be a permutation. The methods below address this inverse
nonnegative representation problem directly.

The permutations are counted as ordered words. Reversal, rotation,
complementation and inversion are not identification rules, and the entries
of $x$ need not be distinct. The problem is an inverse representation
problem: the window sums must be a permutation, whereas the realizing vector
is only required to be nonnegative. For example, $(0,1,2)$ is realized by
$(0,0,0,0,1,1)$, while $(1,0,2)$ is realized by $(1,0,0,0,0,2)$.

A direct scan of all $n!$ permutations becomes expensive even with a linear
feasibility test. The useful feature here is the label zero. If $p_k=0$,
then $x_k=x_{k+1}=x_{k+2}=x_{k+3}=0$. This determines the entire realization
and separates its left and right parts. The enumeration can therefore be
organized by the label sets on those parts, retaining just the last three
realization entries as continuation state.

This paper proves the zero decomposition and its dynamic programming
recurrence, and reports a complete C enumeration through $22$ performed by
the author. After that run, a separately implemented inclusion--exclusion
method recounted all values through $22$. A third method, which generates
from unrestricted initial triples rather than from the zero separator,
checks a smaller, explicitly stated domain. The six values additional to
the historical baseline are given in Section~\ref{sec:results}.
The comparison is with that baseline; no claim of worldwide priority or of
the current remote OEIS endpoint is needed for the numerical results.

\section{Realizations and a feasibility criterion}
\label{sec:realization}

Subtracting consecutive equations in~\eqref{eq:windows} gives
\begin{equation}
 x_{i+4}-x_i=p_{i+1}-p_i\qquad(0\leq i\leq n-2).
 \label{eq:difference}
\end{equation}
Thus the initial four entries determine every subsequent entry.

\begin{lemma}[Uniqueness]
\label{lem:unique}
Every $p\in\calA_n$ has exactly one nonnegative integer realization.
\end{lemma}
\begin{proof}
Let $p_k=0$. Nonnegativity forces $x_k,\ldots,x_{k+3}$ to be zero.
These four consecutive positions include one representative of each residue
class modulo four. Equation~\eqref{eq:difference}, propagated both forward
and backward within each class, determines every other entry. Existence
holds by the definition of $\calA_n$, so the realization is unique.
For $n=1$ the sole window is zero and the same argument applies.
\end{proof}

There is also a direct test that does not require finding the zero position.
For $r\in\{0,1,2,3\}$ put
\[
 K_r=\left\lfloor\frac{n+2-r}{4}\right\rfloor,\qquad
 D_r(0)=0,\qquad
 D_r(k)=\sum_{j=0}^{k-1}(p_{r+4j+1}-p_{r+4j})\quad(1\leq k\leq K_r),
\]
and define
\begin{equation}
 m_r=-\min_{0\leq k\leq K_r}D_r(k).
 \label{eq:minimum}
\end{equation}
All differences appearing here have indices in the range of
\eqref{eq:difference}. In particular $K_r\geq0$, and $m_r\geq0$ because
$D_r(0)=0$.

\begin{proposition}[Residue-minimum criterion]
\label{prop:criterion}
For any word $p\in\Znn^n$, equation~\eqref{eq:windows} has a nonnegative
integer solution if and only if
\begin{equation}
 \sum_{r=0}^3 m_r\leq p_0.
 \label{eq:criterion}
\end{equation}
If $p$ is a permutation of $\{0,\ldots,n-1\}$, feasibility forces equality
in~\eqref{eq:criterion}, and its unique realization is
\begin{equation}
 x_{r+4k}=m_r+D_r(k)\qquad(0\leq r\leq3,\ 0\leq k\leq K_r).
 \label{eq:witness}
\end{equation}
\end{proposition}
\begin{proof}
Write $c_r=x_r$. The difference equations are equivalent to
$x_{r+4k}=c_r+D_r(k)$. These entries are all nonnegative precisely when
$c_r\geq m_r$ for every $r$. The initial window additionally requires
$\sum_r c_r=p_0$. Such nonnegative integer choices exist exactly when
\eqref{eq:criterion} holds: any surplus can be assigned to one $c_r$.
The first window and the consecutive differences then give all the required
windows, proving sufficiency as well as necessity.

For a feasible permutation with $p_k=0$, the realization entries
$x_k,\ldots,x_{k+3}$ are all zero and occupy one position in each residue
class modulo four. The minimum realization entry in each class is therefore
zero, which implies $c_r=m_r$. Summing the first four entries proves equality and
gives~\eqref{eq:witness}.
\end{proof}

Computing the four minima requires $O(n)$ integer operations. This criterion
and reconstruction from the zero position give two useful permutation
membership checks; they are not practical replacements for aggregation over
the full factorial candidate set.

\section{Zero decomposition and weighted subset enumeration}
\label{sec:subset}

\subsection{Outward tails and complementary label sets}

\begin{definition}
\label{def:tail}
An \emph{outward tail} with label word $(v_1,\ldots,v_m)$ starts with
three zero entries. If the current last three entries are $(a,b,c)$, its
next label $v$ appends
\begin{equation}
 d=v-a-b-c\geq0
 \quad\text{and changes the last triple to }(b,c,d).
 \label{eq:transition}
\end{equation}
For a finite set $S$ of positive integers, let $F(S)$ be the number of
orderings of $S$ that are valid outward tails, with $F(\varnothing)=1$.
\end{definition}

For a right tail at a zero window $p_k=0$, the initial triple is
$(x_{k+1},x_{k+2},x_{k+3})=(0,0,0)$; the left tail uses the reversed triple
specified in the proof below. Given the label word, every appended entry is forced,
so $F(S)$ counts words rather than multiple realizations of one word.
The definition uses the actual integer labels, not just their ranks.

\begin{theorem}[Zero decomposition]
\label{thm:convolution}
For $n\geq1$, with $U=\{1,\ldots,n-1\}$, one has
\begin{equation}
 a(n)=\sum_{S\subseteq U}F(S)F(U\setminus S).
 \label{eq:convolution}
\end{equation}
\end{theorem}
\begin{proof}
If the zero of $p$ is at position $k$, its right labels, read outwards,
are $p_{k+1},\ldots,p_{n-1}$. They form a tail starting from
$(x_{k+1},x_{k+2},x_{k+3})=(0,0,0)$. On the left, read the labels
$p_{k-1},\ldots,p_0$ outwards, and the realization entries in reverse
order. This is another tail with initial triple
$(x_{k+2},x_{k+1},x_k)=(0,0,0)$. The label sets of the two tails partition $U$.

Conversely, let $\ell$ and $\rho$ be the label words of tails on $S$ and
$U\setminus S$. Place the entries appended by $\ell$ in reverse order to
the left of four fixed zeros, and those appended by $\rho$ in their generated
order to the right. Both initial triples belong to this same four-zero
block; no window of length four can reach both appended sides.
The resulting window word is
\[
 \bigl(\operatorname{rev}(\ell),\,0,\,\rho\bigr).
\]
Every required window is contained in one of the two tail constructions or
is the central zero window, so this is an admissible permutation. Its zero
position is $|S|$ and recovers both tails uniquely. Hence the construction is
a bijection and the sum counts each ordered permutation exactly once.
The empty-tail case includes $n=1$.
\end{proof}

In general $F(S)$ depends on the values in $S$. For example,
$F(\{1,2,3,4,5\})=2$, with valid orders $(1,2,3,4,5)$ and $(1,2,3,5,4)$,
whereas $F(\{1,2,3,4,7\})=1$. Indeed, the first four labels of any
distinct-label tail must increase strictly; the fifth must be at least the
fourth minus the first. These conditions give the stated examples directly.
Thus replacing $F(S)$ by a function of $|S|$ would change the problem.

\subsection{The weighted recurrence}

Fix a maximum label $N$. Represent $S\subseteq\{1,\ldots,N\}$ by a bit
mask, and let $W(S,a,b,c)$ count the tails on $S$ with last triple $(a,b,c)$.
Initialize $W(\varnothing,0,0,0)=1$, with all other weights zero. For each
unused $v\in\{1,\ldots,N\}\setminus S$ with $v\geq a+b+c$, perform
\begin{equation}
 W(S\cup\{v\},b,c,v-a-b-c)\mathrel{+}=W(S,a,b,c).
 \label{eq:maskedDP}
\end{equation}
Then
\begin{equation}
 F(S)=\sum_{a,b,c\geq0}W(S,a,b,c).
 \label{eq:F}
\end{equation}
Weights, rather than Boolean reachability, are essential: different label
orders may have identical mask and last triple.

\begin{proposition}[Completeness of the recurrence]
\label{prop:DP}
Equations~\eqref{eq:maskedDP}--\eqref{eq:F} compute $F(S)$ exactly for every
subset of $\{1,\ldots,N\}$ and terminate after at most $N$ transitions
along any path.
\end{proposition}
\begin{proof}
The empty tail gives the initial state. Every valid nonempty tail has a
unique final label $v$; removing its final appended entry gives its unique
predecessor tail and the transition~\eqref{eq:maskedDP}. Conversely, every
allowed transition appends a nonnegative integer and an unused label, so
produces a valid tail. Induction on $|S|$ proves that its weight is the
number of histories reaching that full state. Future extensions depend
only on the used labels and the last three entries, so merging equal full
states preserves these multiplicities. Each transition increases $|S|$ by
one, proving termination.
\end{proof}

The triples lie in the finite set
\begin{equation}
 \mathcal T_N=\{(a,b,c)\in\Znn^3:a+b+c\leq N\},
 \qquad |\mathcal T_N|=\binom{N+3}{3}.
 \label{eq:triples}
\end{equation}
This set is closed under~\eqref{eq:transition}: the next triple has sum
$v-a\leq N$. Consequently the number of possible full states is at most
$2^N\binom{N+3}{3}$; only reachable states need storage. A loose bound for
transition attempts is $N2^N\binom{N+3}{3}$, rather than a factorial scan.
This is an upper bound on work, not a prediction of measured running time.

These state and work bounds refer to the unpartitioned recurrence. The
producer stores full mask--triple keys in a hash table with exact equality
comparisons and integer weights, retaining the current and next layers
alongside the $F$ table. It processes disjoint prefix parts separately,
so a full state may be processed in more than one part.

A prefix depth $d$ means fixing the first $d$ tail labels. For $d\leq4$,
these labels must increase strictly: the first four windows are cumulative
sums of the first four appended nonnegative entries, and the labels are
distinct. Conversely any strictly increasing word of at most four positive
labels is realized by its first label and successive positive differences.
Thus $F(S)=1$ for $|S|\leq4$, justifying the initialization of sets shorter
than the chosen depth. Each longer tail belongs to exactly one fixed-prefix
part. The reported run uses $d=1$: nonempty tails are partitioned by their
first label, and the empty tail is added once.

For any $S\subseteq\{1,\ldots,N\}$, a tail on $S$ uses only labels in $S$,
so its construction and $F(S)$ are unchanged by enlarging the ambient label
universe. A table over $\{1,\ldots,N\}$ therefore supplies all coefficients
needed in~\eqref{eq:convolution} for every $n\leq N+1$.

For $n>1$, the summands in~\eqref{eq:convolution} are paired by
$S\leftrightarrow U\setminus S$, with no fixed subset. The implementation
sums masks with their highest bit clear, one member of each pair, and doubles.
This saves work; the answer still
counts all ordered permutations.

\subsection{Bounds on the producer's arithmetic}
\label{sec:prod-arithmetic}

At layer $h$, each label order determines at most one tail, so every state
weight and $F(S)$ with $|S|=h$ is at most $h!$. Even the sum over all states
at that layer is at most $N!/(N-h)!$. Partial sums during a merge or a
partition reduction are nonnegative subcounts of those exact totals.
For $|U|=N$,
\[
 F(S)F(U\setminus S)\leq |S|!(N-|S|)!\leq N!.
\]
The complete convolution is $a(n)\leq n!$ because it counts permutations.
For $n>1$, each partial sum of the half-convolution is at most
$a(n)/2\leq n!/2$, since its summands are nonnegative and the complementary
pairing is exact. Doubling is therefore at most $n!$ as well.
These bounds cover the intermediates, not just the
observed final answer. For the reported run, $N=21$ and
$22!=1124000727777607680000<2^{70}$. The producer uses checked unsigned
128-bit additions and multiplications, through the GCC integer
extension~\cite{gcc}; no floating-point count arithmetic is used.

\section{A second enumeration by inclusion--exclusion}
\label{sec:IE}

The second method permits repeated positive labels and removes them by
requiring full label coverage. Put $N=n-1\geq1$ and $U=\{1,\ldots,N\}$.
For $T\subseteq U$, let $L_h(T)$ count outward tails of length $h$ whose
labels all belong to $T$, with repetition allowed. Here $L_0(T)=1$.
Use the same triple transition~\eqref{eq:transition}, but retain no mask:
\begin{equation}
 \begin{split}
 B_0^T(0,0,0)&=1,\\
 B_{h+1}^T(b,c,v-a-b-c)&\mathrel{+}=B_h^T(a,b,c)
       \quad(v\in T,\ v\geq a+b+c),\\
 L_h(T)&=\sum_{(a,b,c)\in\mathcal T_N} B_h^T(a,b,c).
 \end{split}
 \label{eq:walkDP}
\end{equation}
Other initial weights are zero. Induction on $h$ proves this recurrence
counts exactly the repeated-label tails, since the final step and
predecessor of a tail are unique. It is evaluated only for $0\leq h\leq N$.
Define
\begin{equation}
 Q(T)=\sum_{k=0}^N L_k(T)L_{N-k}(T).
 \label{eq:Q}
\end{equation}
The index $k$ is the position of the zero window. Left tails are again read
outwards and then reversed to form the window word. Every other window has
a positive label, so the zero is unique even though positive labels can repeat.

\begin{theorem}[Inclusion--exclusion count]
\label{thm:IE}
For $n\geq2$,
\begin{equation}
 a(n)=\sum_{T\subseteq U}(-1)^{N-|T|}Q(T).
 \label{eq:IE}
\end{equation}
\end{theorem}
\begin{proof}
Fix a pair of tails of total length $N$ contributing to~\eqref{eq:Q};
its left length fixes $k$. Let $V\subseteq U$
be the union of the labels they use. Its total coefficient in the right
side of~\eqref{eq:IE} is
\[
 \sum_{V\subseteq T\subseteq U}(-1)^{N-|T|}
 =\begin{cases}1,&V=U,\\0,&V\ne U.\end{cases}
\]
For $V\ne U$, this follows by expanding $(1-1)^{N-|V|}$; for $V=U$ there
is one summand. This is Boolean-lattice M\"obius inversion, whose
inclusion--exclusion sign is described in~\cite[Section 3]{rota}.
There are exactly $N$ positive window positions. Covering all $N$ positive
labels therefore forces each to occur exactly once: a repetition would leave
fewer than $N$ distinct labels in those $N$ positions. The surviving pairs
are precisely those of Theorem~\ref{thm:convolution}, with coefficient one.
\end{proof}

The C++17 checker visits every one of the $2^N$ subsets, using dense arrays
on $\mathcal T_N$. Unlike the producer, it uses neither distinct-label
mask states within a walk nor complementary $F$ tables, prefix partitions,
hash-table aggregation, or a half-convolution. Its only input is $n$;
the checker reads and compares no expected sequence count. Agreement with
the native table is established externally from the completed output records.
A direct
implementation has an upper work bound
$O(2^N N^2\binom{N+3}{3})$, including all $N$ layers and at most $N$ label
choices per state. The count arrays use $O(\binom{N+3}{3}+N)$ storage.
The implementation also precomputes $O(N\binom{N+3}{3})$ transition indices
and an $O(N^3)$ triple-index array, giving total auxiliary storage
$O(N\binom{N+3}{3}+N^3)$.

\subsection{Signed cancellation and intermediate bounds}
\label{sec:IE-arithmetic}

Each tail step has at most $N$ choices. For $0\leq h\leq N$,
\begin{equation}
 L_h(T)\leq N^h,\qquad Q(T)\leq(N+1)N^N.
 \label{eq:walkbound}
\end{equation}
Every state weight, layer partial total and product $L_k(T)L_{N-k}(T)$ is
bounded by the corresponding nonnegative total above. In any order of
subset processing, the absolute value of a signed partial sum
in~\eqref{eq:IE} is at most
\begin{equation}
 2^N(N+1)N^N.
 \label{eq:absolute}
\end{equation}
At $n=22$, this is
\[
 269561447117208400366948292518477824<2^{118}<2^{127}.
\]
The bound increases with $N\geq1$, so it covers all reported cases
$2\leq n\leq22$. Unsigned 128-bit integers store nonnegative walk counts
and products; a signed 128-bit integer stores the cancelling sum. The
checker verifies layer bounds, per-subset bounds, subset coverage and the
final interval $[0,n!]$. The case $n=1$ is the single all-zero window and
is counted as one. This arithmetic argument is the one applying to the
checker used for the reported $22$-term validation.

\section{Structural consequences}
\label{sec:structure}

\begin{proposition}[Reversal and parity]
\label{prop:parity}
For $n>1$, $a(n)$ is even.
\end{proposition}
\begin{proof}
Reversing a realization $x$ reverses its window word, so reversal preserves
$\calA_n$. A fixed word under reversal would have $p_0=p_{n-1}$, impossible
for a permutation of length greater than one. Thus reversal pairs the
admissible permutations without fixed points.
\end{proof}

The following observation was formulated after the numerical extension,
during preparation for the manuscript. It is not used as a pruning rule in
the enumeration.

\begin{proposition}[Extension at either end]
\label{prop:growth}
For every $n\geq1$,
\begin{equation}
 a(n+1)\geq2a(n).
 \label{eq:growth}
\end{equation}
\end{proposition}
\begin{proof}
Let $p\in\calA_n$ have realization $x$. Set
$s=x_n+x_{n+1}+x_{n+2}$. Its last old window is $x_{n-1}+s\leq n-1$,
so $s\leq n-1$. Append the nonnegative entry $n-s$ to $x$. This preserves
all old windows and creates the final window $n$, giving $(p,n)\in\calA_{n+1}$.

Similarly let $t=x_0+x_1+x_2\leq p_0\leq n-1$. Prepending $n-t$ to $x$
gives $(n,p)\in\calA_{n+1}$. Each map is injective because deleting the
added label recovers $p$. Their images are disjoint: the unique maximum $n$
is at opposite ends, which are distinct positions even when $n=1$.
Together the maps prove~\eqref{eq:growth}.
\end{proof}

\begin{corollary}
\label{cor:bounds}
For every $n\geq1$,
\[
 2^{n-1}\leq a(n)\leq n!,
\]
and $a(n)/2^{n-1}$ is nondecreasing.
\end{corollary}
\begin{proof}
The lower bound follows from $a(1)=1$ and Proposition~\ref{prop:growth}.
The upper bound counts all permutations. Dividing~\eqref{eq:growth} by
$2^n$ gives the monotonicity.
\end{proof}

These are uniform mathematical statements. They do not determine an exact
recurrence or an asymptotic growth rate from the finite table. Parity and
the growth inequality are consistency checks, not independent exhaustive
enumerations.

\section{Computation and validation}
\label{sec:validation}

\subsection{Implementation and observed resources}

The frozen producer is written in C17 with OpenMP and the GCC 128-bit
extension. The completed run used maximum label $21$ ($n=22$), one worker,
and prefix depth one. It computes the whole prefix $1\leq n\leq22$ by
reusing the tail table for the smaller label sets. The source, build recipe,
exact run records and validator records are provided with the computational
package~\cite{corti}.

The recorded host was an AMD Ryzen 9 7940HS with eight cores and sixteen
hardware threads, running Linux Mint 22.3 on x86-64, with recorded kernel
6.17.0-42-generic. The producer
was built with GCC 13.3.0, C17, \texttt{-O3}, \texttt{-march=znver4},
\texttt{-mtune=znver4}, link-time optimization and \texttt{-fopenmp}.
Its build also enables loop unrolling, frame-pointer omission and warning
checks, and links the system cryptographic library for checkpoint integrity.
The C++17 checker used GCC 13.3.0 with \texttt{-O3} and the same architecture
tuning, without OpenMP. The unrestricted-window checker used Python 3.12.3
and arbitrary-precision integers. Each reported count was performed by a
single worker or serial process; the processor's core count is not a claim
of parallel execution.

The native run took about $297$ seconds wall time, processing
$1177119126$ state occurrences and $2911882132$ transitions
across its prefix parts. Its peak tracked allocations were $3486204032$
bytes, approximately $3.25$ GiB; this is not a measured process RSS.
The inclusion--exclusion count at $22$ took about $251$ seconds and had a
measured peak RSS of $16404$ KiB, approximately $16$ MiB. Its dense triple
state set has $\binom{24}{3}=2024$ elements. The unrestricted Python count
at $17$ took about $102$ seconds and had a measured peak RSS of
$4334392$ KiB, approximately $4.13$ GiB. These are observations for the
stated configurations, not projections to larger indices.

\subsection{Different checks and their coverage}

Table~\ref{tab:validation} records the complete domains actually checked.
For the direct Python method, start with every nonnegative triple
$(a,b,c)$ satisfying $a+b+c\leq n-1$, each of weight one, and an empty
used-label mask. Apply~\eqref{eq:transition} for unused labels
$v\in\{0,\ldots,n-1\}$ for exactly $n$ steps. Every witness is included:
its first triple has sum at most its first window, and its successive
windows give an allowed path. Conversely every complete path is a
nonnegative witness with all $n$ labels used. Lemma~\ref{lem:unique}
then makes the path count a permutation count. This method does not
generate by a zero separator or combine complementary tails.

The complete Python results at $1,\ldots,12$ and $16$ were obtained with
earlier inline implementations of this recurrence, retained with their
original commands and outputs. The retained parameterized program was
qualified at $1,2,6$ and produced the complete result at $17$.
The combined row below describes the method's recorded coverage across
these implementations, not execution of one frozen file at every listed index.

\begin{table}[htbp]
\centering
\small
\begin{tabular}{@{}>{\raggedright\arraybackslash}p{0.33\textwidth}>{\raggedright\arraybackslash}p{0.19\textwidth}>{\raggedright\arraybackslash}p{0.40\textwidth}@{}}
\toprule
Method & Complete indices & Role or comparison \\
\midrule
C weighted subset enumeration & $1,\ldots,22$ & Native complete counts \\
C++ repeated-label walks and inclusion--exclusion & $1,\ldots,22$ & Full recount by a distinct recurrence and implementation \\
Python unrestricted initial triples & $1,\ldots,12,16,17$ & Distinct generation without zero decomposition, using arbitrary-precision counts \\
All permutations: residue test and reconstruction from zero & $1,\ldots,9$ & Exhaustive comparison of membership predicates and totals \\
\bottomrule
\end{tabular}
\caption{Recorded complete domains. Comparisons with the native counts agree.
The Python row combines the retained inline and parameterized implementations and
does not include $13,14,15$ or any index above $17$.}
\label{tab:validation}
\end{table}

Before checking the six additional rows, inclusion--exclusion reproduced the
historical $1,\ldots,16$ prefix. Each run at $17,\ldots,22$ then visited
all $2^{n-1}$ subsets, including $2097152$ subsets at $22$, and matched
the corresponding native result. The first additional value, at $17$,
was also fully recomputed by the unrestricted Python method.

The producer computes its values before an output check compares indices
$1,\ldots,16$ against an embedded historical table and rejects completion
on disagreement. That check supplies no count to the recurrence and contains
no expected value for $17,\ldots,22$. The IE checker has no expected-count
table; its results are compared with the native values externally.

Earlier Python tail-DP tests reproduced the known prefixes with maximum
index $12,13,14,15,16$. At maximum index $16$, the complete table of
$32768$ coefficients $F(S)$ agreed across prefix depths $0,1,2,3,4$.
Those tests support implementation consistency;
they are not additional independent mathematical models.

The producer and the inclusion--exclusion checker share the zero-anchor
theorem, but their aggregation and repeated-label cancellation are different.
The unrestricted method checks that shared generation assumption on its
stated complete domain. C and C++ also share a compiler family; arbitrary
precision in Python supplies a different arithmetic check at the reported
indices. All local validation belongs to the same OpenAI/Codex maintenance
context. Here ``independent'' describes the indicated algorithmic or
implementation separation, not a third-party replication of the computations.

Original runnable checkers, inputs and completion records are included
in the computational package~\cite{corti}. File-identity and result-consistency checks reuse that evidence on
unchanged code; they do not constitute fresh enumerations. Incomplete runs,
interrupted processes and partial state counts contribute no sequence term.

\section{Results}
\label{sec:results}

\begin{theorem}[Computational result]
\label{thm:result}
The complete enumerations described in Sections~\ref{sec:subset}
and~\ref{sec:IE} agree for every $1\leq n\leq22$ and give the values
in Table~\ref{tab:counts}. In particular,
\[
 a(22)=2055385215352.
\]
\end{theorem}
\begin{proof}[Computational proof]
The masked recurrence is exhaustive by Proposition~\ref{prop:DP}, and
Theorem~\ref{thm:convolution} converts its tail counts to the required
ordered permutation counts. The completed native run produced all $22$
rows. The second enumeration uses the exhaustive walk recurrence
\eqref{eq:walkDP} and Theorem~\ref{thm:IE}, with complete subset coverage
at each index. The retained completion records and exact result comparisons
agree row for row. The bounds in Sections~\ref{sec:prod-arithmetic}
and~\ref{sec:IE-arithmetic} cover the fixed-width intermediates.
The additional checks and their limits are those in
Table~\ref{tab:validation}. This is a finite computational result supported
by the recorded executions, not a closed formula beyond their range.
\end{proof}

\begin{table}[htbp]
\centering
\begin{tabular}{@{}rr@{\qquad}rr@{}}
\toprule
$n$ & $a(n)$ & $n$ & $a(n)$ \\
\midrule
1 & 1 & 12 & 67646 \\
2 & 2 & 13 & 301210 \\
3 & 4 & 14 & 1397898 \\
4 & 8 & 15 & 6786792 \\
5 & 16 & 16 & 35154332 \\
6 & 34 & 17 & 193009094 \\
7 & 82 & 18 & 1125920768 \\
8 & 252 & 19 & 6895579278 \\
9 & 890 & 20 & 44276337188 \\
10 & 3644 & 21 & 295719612912 \\
11 & 15390 & 22 & 2055385215352 \\
\bottomrule
\end{tabular}
\caption{The validated contiguous prefix. Rows $1$--$16$ reproduce the
historical OEIS baseline attributed to Hardin~\cite{hardin}; rows
$17$--$22$ are the six locally added counts, each fully recounted by
inclusion--exclusion.}
\label{tab:counts}
\end{table}

The canonical local b-file contains these $22$ indexed values and preserves
the historical prefix. Every row satisfies the proved parity and growth
constraints. The first index without a complete validated value in this
study is $23$. This endpoint describes the completed and validated
computations, rather than a mathematical limit of the methods. Neither
those constraints nor the tail identities evaluate
it from the existing table without further enumeration. No count at $23$
or beyond is asserted here.

\section{Conclusion}

The zero label converts an inverse nonnegative-window problem into an exact
pairing of tails on complementary positive label sets. It both removes
realization multiplicity and permits finite-state aggregation without scanning
all permutations. Masked distinct-label enumeration and repeated-label
inclusion--exclusion give matching complete counts through $22$, extending
the study's historical baseline by six values. The residue-minimum criterion
decides realizability directly, while reversal and the two end-extension
injections give evenness and a uniform exponential lower bound. The finite counts
do not determine an asymptotic law; the useful outcome is an explicitly
proved counting framework and a validated contiguous numerical prefix.

\section*{Data availability}

The producer, build recipe, full count table, validation programs, original
inputs and execution records are available in the GitHub repository and
the Zenodo release series~\cite{corti}. The manuscript source and its
matching PDF are included in the computational package. The cited Zenodo
DOI identifies the release series; individual deposited releases have
distinct version-specific DOIs.

\section*{AI use disclosure}

Large language models were used as research, code-review and drafting
assistants during the A180207 workflow. Tool and vendor names below follow
the supplied records; exact model versions and backends were not consistently
authenticated.

\noindent\textbf{Mathematical analysis.}
OpenAI/Codex assisted with the window-difference analysis, zero decomposition,
feasibility and arithmetic proofs, and the later end-extension observation.
NotebookLM supported consultation of prior methodological notes during
strategy preparation; the sequence-specific identities were derived and
checked locally.
The mathematical claims are supported by the explicit arguments, primary
sources and retained computational artifacts, not by conversational assertions
or agreement among language models.

\smallskip
\noindent\textbf{Algorithm and code.}
OpenAI/Codex assisted with producer development, maintenance, validation design
and local evidence checks. The author performed the completed native run.
The IE and Python validation processes were executed locally through
OpenAI/Codex under the author's authorized validation scope.
Received code-review reports identified their tools as ChatGPT, Claude,
DeepSeek, GLM, Grok, Meta, MiniMax and Qwen; further reports came from Claude,
Meta and MiniMax for follow-up code checks. These names describe the tool
or vendor designations in the supplied
reports; backend identities and reported remote executions were not
independently authenticated. Exact model versions were not consistently
verified. Review opinions are distinct from the local execution evidence,
the complete inclusion--exclusion recounts and the stated Python checks.

\smallskip
\noindent\textbf{Manuscript.}
OpenAI/Codex assisted with drafting, source checking and internal review of
this manuscript. Numerical and validation claims were checked against the
b-file, native result table, original completion records and validator
sources. This internal work is not external manuscript peer review.
Manuscript-review reports supplied by the author bore the designations
ChatGPT, Claude, DeepSeek, Gemini, GLM, Grok, Meta, MiniMax and Qwen.
These reports were evaluated critically; their designations do not
authenticate distinct model families, and reported external computations
were not added to the verified local validation domains without executable
evidence. They are separate from the internal drafting checks.

\smallskip
The author, Carlo Corti, takes full scientific responsibility for the
mathematical statements, algorithm, implementation, numerical results,
validation boundaries and conclusions of this paper.

\clearpage
\begin{thebibliography}{9}
\raggedright

\bibitem{hardin}
R. H. Hardin,
\emph{A180207: Number of permutations of $0,\ldots,n-1$ representable as
consecutive sums of 4 adjacent elements of a sequence of $n+3$ nonnegative
integers}, The On-Line Encyclopedia of Integer Sequences,
OEIS author date 15 August 2010.
\url{https://oeis.org/A180207}.
Historical $16$-term entry snapshot used in this study; consulted
November 2026.

\bibitem{family}
The OEIS Foundation Inc.,
\emph{Consecutive-sum window-length variants}, entries
A180205, A180206, A180208, A180209 and A180210,
The On-Line Encyclopedia of Integer Sequences.
Definitions listed with the A180207 entry; consulted November 2026.
\url{https://oeis.org/A180205},
\url{https://oeis.org/A180206},
\url{https://oeis.org/A180208},
\url{https://oeis.org/A180209},
\url{https://oeis.org/A180210}.

\bibitem{konieczny}
J. Konieczny,
\emph{On consecutive sums in permutations},
arXiv:1504.07156v5, 27 August 2021.
Preprint DOI: 10.48550/arXiv.1504.07156.
\url{https://arxiv.org/abs/1504.07156v5}.

\bibitem{higashitani}
A. Higashitani and K. Kurimoto,
\emph{Permutations with small maximal $k$-consecutive sums},
arXiv:1801.00416v4, 27 May 2019.
Preprint DOI: 10.48550/arXiv.1801.00416.
\url{https://arxiv.org/abs/1801.00416v4}.

\bibitem{rota}
G.-C. Rota,
On the foundations of combinatorial theory I. Theory of M\"obius functions,
\emph{Zeitschrift f\"ur Wahrscheinlichkeitstheorie und Verwandte Gebiete}
\textbf{2} (1964), 340--368.
\href{https://link.springer.com/article/10.1007/BF00531932}{doi:10.1007/BF00531932}.

\bibitem{gcc}
Free Software Foundation,
\emph{Using the GNU Compiler Collection (GCC)}, subsection ``128-bit integers.''
\url{https://gcc.gnu.org/onlinedocs/gcc/_005f_005fint128.html}.
Accessed November 2026.

\bibitem{corti}
C. Corti,
\emph{Exact Enumeration of Four-Window Permutations: OEIS A180207 through
$n=22$}, GitHub repository and Zenodo archive, 2026.
\url{https://github.com/carcorti/A180207};
\url{https://doi.org/10.5281/zenodo.24194032}.
The Zenodo DOI identifies the release series.

\end{thebibliography}
\end{document}
